\IfFileExists{stacks-project.cls}{%
\documentclass{stacks-project}
}{%
\documentclass{amsart}
}

% For dealing with references we use the comment environment
\usepackage{verbatim}
\newenvironment{reference}{\comment}{\endcomment}
%\newenvironment{reference}{}{}
\newenvironment{slogan}{\comment}{\endcomment}
\newenvironment{history}{\comment}{\endcomment}

% For commutative diagrams we use Xy-pic
\usepackage[all]{xy}

% We use 2cell for 2-commutative diagrams.
\xyoption{2cell}
\UseAllTwocells

% We use multicol for the list of chapters between chapters
\usepackage{multicol}

% This is generally recommended for better output
\usepackage{lmodern}
\usepackage[T1]{fontenc}

% For cross-file-references
\usepackage{xr-hyper}

% Package for hypertext links:
\usepackage{hyperref}

% For any local file, say "hello.tex" you want to link to please
% use \externaldocument[hello-]{hello}
\externaldocument[introduction-]{introduction}
\externaldocument[conventions-]{conventions}
\externaldocument[sets-]{sets}
\externaldocument[categories-]{categories}
\externaldocument[topology-]{topology}
\externaldocument[sheaves-]{sheaves}
\externaldocument[sites-]{sites}
\externaldocument[stacks-]{stacks}
\externaldocument[fields-]{fields}
\externaldocument[algebra-]{algebra}
\externaldocument[brauer-]{brauer}
\externaldocument[homology-]{homology}
\externaldocument[derived-]{derived}
\externaldocument[simplicial-]{simplicial}
\externaldocument[more-algebra-]{more-algebra}
\externaldocument[smoothing-]{smoothing}
\externaldocument[modules-]{modules}
\externaldocument[sites-modules-]{sites-modules}
\externaldocument[injectives-]{injectives}
\externaldocument[cohomology-]{cohomology}
\externaldocument[sites-cohomology-]{sites-cohomology}
\externaldocument[dga-]{dga}
\externaldocument[dpa-]{dpa}
\externaldocument[sdga-]{sdga}
\externaldocument[hypercovering-]{hypercovering}
\externaldocument[schemes-]{schemes}
\externaldocument[constructions-]{constructions}
\externaldocument[properties-]{properties}
\externaldocument[morphisms-]{morphisms}
\externaldocument[coherent-]{coherent}
\externaldocument[divisors-]{divisors}
\externaldocument[limits-]{limits}
\externaldocument[varieties-]{varieties}
\externaldocument[topologies-]{topologies}
\externaldocument[descent-]{descent}
\externaldocument[perfect-]{perfect}
\externaldocument[more-morphisms-]{more-morphisms}
\externaldocument[flat-]{flat}
\externaldocument[groupoids-]{groupoids}
\externaldocument[more-groupoids-]{more-groupoids}
\externaldocument[etale-]{etale}
\externaldocument[chow-]{chow}
\externaldocument[intersection-]{intersection}
\externaldocument[pic-]{pic}
\externaldocument[weil-]{weil}
\externaldocument[adequate-]{adequate}
\externaldocument[dualizing-]{dualizing}
\externaldocument[duality-]{duality}
\externaldocument[discriminant-]{discriminant}
\externaldocument[derham-]{derham}
\externaldocument[local-cohomology-]{local-cohomology}
\externaldocument[algebraization-]{algebraization}
\externaldocument[curves-]{curves}
\externaldocument[resolve-]{resolve}
\externaldocument[models-]{models}
\externaldocument[functors-]{functors}
\externaldocument[equiv-]{equiv}
\externaldocument[pione-]{pione}
\externaldocument[etale-cohomology-]{etale-cohomology}
\externaldocument[proetale-]{proetale}
\externaldocument[relative-cycles-]{relative-cycles}
\externaldocument[more-etale-]{more-etale}
\externaldocument[trace-]{trace}
\externaldocument[crystalline-]{crystalline}
\externaldocument[spaces-]{spaces}
\externaldocument[spaces-properties-]{spaces-properties}
\externaldocument[spaces-morphisms-]{spaces-morphisms}
\externaldocument[decent-spaces-]{decent-spaces}
\externaldocument[spaces-cohomology-]{spaces-cohomology}
\externaldocument[spaces-limits-]{spaces-limits}
\externaldocument[spaces-divisors-]{spaces-divisors}
\externaldocument[spaces-over-fields-]{spaces-over-fields}
\externaldocument[spaces-topologies-]{spaces-topologies}
\externaldocument[spaces-descent-]{spaces-descent}
\externaldocument[spaces-perfect-]{spaces-perfect}
\externaldocument[spaces-more-morphisms-]{spaces-more-morphisms}
\externaldocument[spaces-flat-]{spaces-flat}
\externaldocument[spaces-groupoids-]{spaces-groupoids}
\externaldocument[spaces-more-groupoids-]{spaces-more-groupoids}
\externaldocument[bootstrap-]{bootstrap}
\externaldocument[spaces-pushouts-]{spaces-pushouts}
\externaldocument[spaces-chow-]{spaces-chow}
\externaldocument[groupoids-quotients-]{groupoids-quotients}
\externaldocument[spaces-more-cohomology-]{spaces-more-cohomology}
\externaldocument[spaces-simplicial-]{spaces-simplicial}
\externaldocument[spaces-duality-]{spaces-duality}
\externaldocument[formal-spaces-]{formal-spaces}
\externaldocument[restricted-]{restricted}
\externaldocument[spaces-resolve-]{spaces-resolve}
\externaldocument[formal-defos-]{formal-defos}
\externaldocument[defos-]{defos}
\externaldocument[cotangent-]{cotangent}
\externaldocument[examples-defos-]{examples-defos}
\externaldocument[algebraic-]{algebraic}
\externaldocument[examples-stacks-]{examples-stacks}
\externaldocument[stacks-sheaves-]{stacks-sheaves}
\externaldocument[criteria-]{criteria}
\externaldocument[artin-]{artin}
\externaldocument[quot-]{quot}
\externaldocument[stacks-properties-]{stacks-properties}
\externaldocument[stacks-morphisms-]{stacks-morphisms}
\externaldocument[stacks-limits-]{stacks-limits}
\externaldocument[stacks-cohomology-]{stacks-cohomology}
\externaldocument[stacks-perfect-]{stacks-perfect}
\externaldocument[stacks-introduction-]{stacks-introduction}
\externaldocument[stacks-more-morphisms-]{stacks-more-morphisms}
\externaldocument[stacks-geometry-]{stacks-geometry}
\externaldocument[moduli-]{moduli}
\externaldocument[moduli-curves-]{moduli-curves}
\externaldocument[examples-]{examples}
\externaldocument[exercises-]{exercises}
\externaldocument[guide-]{guide}
\externaldocument[desirables-]{desirables}
\externaldocument[coding-]{coding}
\externaldocument[obsolete-]{obsolete}
\externaldocument[fdl-]{fdl}
\externaldocument[index-]{index}

% Theorem environments.
%
\theoremstyle{plain}
\newtheorem{theorem}[subsection]{Theorem}
\newtheorem{proposition}[subsection]{Proposition}
\newtheorem{lemma}[subsection]{Lemma}

\theoremstyle{definition}
\newtheorem{definition}[subsection]{Definition}
\newtheorem{example}[subsection]{Example}
\newtheorem{exercise}[subsection]{Exercise}
\newtheorem{situation}[subsection]{Situation}

\theoremstyle{remark}
\newtheorem{remark}[subsection]{Remark}
\newtheorem{remarks}[subsection]{Remarks}

\numberwithin{equation}{subsection}

% Macros
%
\def\lim{\mathop{\mathrm{lim}}\nolimits}
\def\colim{\mathop{\mathrm{colim}}\nolimits}
\def\Spec{\mathop{\mathrm{Spec}}}
\def\Hom{\mathop{\mathrm{Hom}}\nolimits}
\def\Ext{\mathop{\mathrm{Ext}}\nolimits}
\def\SheafHom{\mathop{\mathcal{H}\!\mathit{om}}\nolimits}
\def\SheafExt{\mathop{\mathcal{E}\!\mathit{xt}}\nolimits}
\def\Sch{\mathit{Sch}}
\def\Mor{\mathop{\mathrm{Mor}}\nolimits}
\def\Ob{\mathop{\mathrm{Ob}}\nolimits}
\def\Sh{\mathop{\mathit{Sh}}\nolimits}
\def\NL{\mathop{N\!L}\nolimits}
\def\CH{\mathop{\mathrm{CH}}\nolimits}
\def\proetale{{pro\text{-}\acute{e}tale}}
\def\etale{{\acute{e}tale}}
\def\QCoh{\mathit{QCoh}}
\def\Ker{\mathop{\mathrm{Ker}}}
\def\Im{\mathop{\mathrm{Im}}}
\def\Coker{\mathop{\mathrm{Coker}}}
\def\Coim{\mathop{\mathrm{Coim}}}

% Boxtimes
%
\DeclareMathSymbol{\boxtimes}{\mathbin}{AMSa}{"02}

%
% Macros for moduli stacks/spaces
%
\def\QCohstack{\mathcal{QC}\!\mathit{oh}}
\def\Cohstack{\mathcal{C}\!\mathit{oh}}
\def\Spacesstack{\mathcal{S}\!\mathit{paces}}
\def\Quotfunctor{\mathrm{Quot}}
\def\Hilbfunctor{\mathrm{Hilb}}
\def\Curvesstack{\mathcal{C}\!\mathit{urves}}
\def\Polarizedstack{\mathcal{P}\!\mathit{olarized}}
\def\Complexesstack{\mathcal{C}\!\mathit{omplexes}}
% \Pic is the operator that assigns to X its picard group, usage \Pic(X)
% \Picardstack_{X/B} denotes the Picard stack of X over B
% \Picardfunctor_{X/B} denotes the Picard functor of X over B
\def\Pic{\mathop{\mathrm{Pic}}\nolimits}
\def\Picardstack{\mathcal{P}\!\mathit{ic}}
\def\Picardfunctor{\mathrm{Pic}}
\def\Deformationcategory{\mathcal{D}\!\mathit{ef}}
\usepackage{cleveref}
\crefname{originallemma}{original lemma}{original lemmas}
\Crefname{originallemma}{Original lemma}{Original lemmas}
\makeatletter
\newlabel{schemes-lemma-morphism-into-affine}{{01I1}{source}{}{sourcepointer-1}{}}
\newlabel{schemes-lemma-morphism-into-affine@cref}{{[originallemma][0][]01I1}{[1][0][]source}{}{}{}}
\makeatother
\crefname{originallemma}{original lemma}{original lemmas}
\Crefname{originallemma}{Original lemma}{Original lemmas}
\makeatletter
\newlabel{limits-lemma-topology-limit}{{01YY}{source}{}{sourcepointer-2}{}}
\newlabel{limits-lemma-topology-limit@cref}{{[originallemma][0][]01YY}{[1][0][]source}{}{}{}}
\makeatother
\crefname{originallemma}{original lemma}{original lemmas}
\Crefname{originallemma}{Original lemma}{Original lemmas}
\makeatletter
\newlabel{topology-lemma-make-spectral-space}{{0A31}{source}{}{sourcepointer-3}{}}
\newlabel{topology-lemma-make-spectral-space@cref}{{[originallemma][0][]0A31}{[1][0][]source}{}{}{}}
\makeatother
\crefname{originallemma}{original lemma}{original lemmas}
\Crefname{originallemma}{Original lemma}{Original lemmas}
\makeatletter
\newlabel{properties-lemma-cartesian-diagram-quasi-affine}{{0ARY}{source}{}{sourcepointer-4}{}}
\newlabel{properties-lemma-cartesian-diagram-quasi-affine@cref}{{[originallemma][0][]0ARY}{[1][0][]source}{}{}{}}
\makeatother
\title{An intersection of quasi-compact opens in an ind-quasi-affine scheme identifies with an intersection of quasi-compact opens in an affine scheme, and opens of X inside E map to opens of Y}
\author{The Stacks Project Authors}
\date{stacks-project:more-groupoids.tex:local-snapshot-20261012; GNU FDL 1.2 or later}
\begin{document}
\maketitle
\setcounter{subsection}{1}
\noindent\textit{Editorial note.} All delivered mathematical text is the author's own, copied from the pinned Stacks Project chapter file more-groupoids.tex; no statement or proof step was rewritten, reordered or removed. The counted span is the author environment and its own proof environment. References inside the span are handed to the assembler as navigation stubs that carry the canonical tag of the cited tag only; whether each of them is load-bearing is recorded on the stub itself.\par\medskip
\noindent\textit{Allowed standard prerequisites (editorial).}\begin{itemize}
\item The canonical morphism from a locally ringed space with a section-generating global sections map into the spectrum of its ring of sections (schemes chapter, Lemma schemes-lemma-morphism-into-affine, canonical tag 01I1). cited at statement lines 3215-3219 and proof lines 3227-3229. Produces the canonical morphism j : (E, O\_X|\_E) to (Spec(A), O\_Spec(A)) whose image identification is the content of the lemma.
\item Anti-equivalence of colimits of rings and limits of affine schemes on underlying topological spaces (limits chapter, Lemma limits-lemma-topology-limit, canonical tag 01YY). cited at proof lines 3244-3245. From A = colim A\_i it gives Spec(A) = lim Spec(A\_i) as topological spaces, the limit through which the image E' is computed.
\item Presentation of a spectral-like intersection as a limit of schemes (topology chapter, Lemma topology-lemma-make-spectral-space, canonical tag 0A31). cited at proof lines 3246-3247. Gives E = lim U\_i, which combined with the limit identification shows E maps homeomorphically onto its image E' inside Spec(A).
\item Cartesianity of the square attached to a quasi-compact open of a quasi-affine scheme (properties chapter, Lemma properties-lemma-cartesian-diagram-quasi-affine, canonical tag 0ARY). cited at proof lines 3273-3275. Makes the first of the three squares cartesian, hence the second, and after passing to the limit the third, whose top horizontal arrow is then an open immersion.
\end{itemize}\medskip
\ifdefined\hypertarget\hypertarget{sourcepointer-1}{}\fi
\noindent\textit{Source reference.} Original lemma 01I1: sources/stacks-chow/tags.txt, source lines 1962--1962. Cited at statement lines 3215-3219 for the canonical morphism j and at proof lines 3227-3229 for its applicability to the locally ringed space (E, O\_X|\_E) with A = Gamma(E, O\_X|\_E).\par\medskip
\ifdefined\hypertarget\hypertarget{sourcepointer-2}{}\fi
\noindent\textit{Source reference.} Original lemma 01YY: sources/stacks-chow/tags.txt, source lines 2573--2573. Cited at proof lines 3244-3245: A = colim A\_i gives Spec(A) = lim Spec(A\_i) as topological spaces.\par\medskip
\ifdefined\hypertarget\hypertarget{sourcepointer-3}{}\fi
\noindent\textit{Source reference.} Original lemma 0A31: sources/stacks-chow/tags.txt, source lines 12724--12724. Cited at proof lines 3246-3247: E = lim U\_i, the input for the homeomorphism of E onto its image in Spec(A).\par\medskip
\ifdefined\hypertarget\hypertarget{sourcepointer-4}{}\fi
\noindent\textit{Source reference.} Original lemma 0ARY: sources/stacks-chow/tags.txt, source lines 13561--13561. Cited at proof lines 3273-3275: the first square is cartesian, hence the second, and the limit square whose top arrow is an open immersion.\par\medskip
\begin{lemma}
\label{lemma-sits-in-functions}
Let $X$ be an ind-quasi-affine scheme. Let $E \subset X$ be an
intersection of a nonempty family of quasi-compact opens of $X$.
Set $A = \Gamma(E, \mathcal{O}_X|_E)$ and $Y = \Spec(A)$.
Then the canonical morphism
$$
j : (E, \mathcal{O}_X|_E) \longrightarrow (Y, \mathcal{O}_Y)
$$
of Schemes, Lemma \ref{schemes-lemma-morphism-into-affine}
determines an isomorphism
$(E, \mathcal{O}_X|_E) \to (E', \mathcal{O}_Y|_{E'})$
where $E' \subset Y$ is an intersection of quasi-compact opens.
If $W \subset E$ is open in $X$, then $j(W)$ is open in $Y$.
\end{lemma}

\begin{proof}
Note that $(E, \mathcal{O}_X|_E)$ is a locally ringed space so that
Schemes, Lemma \ref{schemes-lemma-morphism-into-affine} applies
to $A \to \Gamma(E, \mathcal{O}_X|_E)$. Write $E = \bigcap_{i \in I} U_i$
with $I \not = \emptyset$ and $U_i \subset X$ quasi-compact open.
We may and do assume that for $i, j \in I$ there exists a $k \in I$ with
$U_k \subset U_i \cap U_j$. Set $A_i = \Gamma(U_i, \mathcal{O}_{U_i})$.
We obtain commutative diagrams
$$
\xymatrix{
(E, \mathcal{O}_X|_E) \ar[r] \ar[d] &
(\Spec(A), \mathcal{O}_{\Spec(A)}) \ar[d] \\
(U_i, \mathcal{O}_{U_i}) \ar[r] &
(\Spec(A_i), \mathcal{O}_{\Spec(A_i)})
}
$$
Since $U_i$ is quasi-affine, we see that $U_i \to \Spec(A_i)$
is a quasi-compact open immersion. On the other hand
$A = \colim A_i$. Hence $\Spec(A) = \lim \Spec(A_i)$ as topological
spaces (Limits, Lemma \ref{limits-lemma-topology-limit}). Since
$E = \lim U_i$ (by Topology, Lemma \ref{topology-lemma-make-spectral-space})
we see that $E \to \Spec(A)$ is a homeomorphism onto its
image $E'$ and that $E'$ is the intersection of the inverse images
of the opens $U_i \subset \Spec(A_i)$ in $\Spec(A)$. For any
$e \in E$ the local ring $\mathcal{O}_{X, e}$ is the value
of $\mathcal{O}_{U_i, e}$ which is the same as the value on $\Spec(A)$.

\medskip\noindent
To prove the final assertion of the lemma we argue as follows.
Pick $i, j \in I$ with $U_i \subset U_j$.
Consider the following commutative diagrams
$$
\xymatrix{
U_i \ar[r] \ar[d] & \Spec(A_i) \ar[d] \\
U_i \ar[r] & \Spec(A_j)
}
\quad\quad
\xymatrix{
W \ar[r] \ar[d] & \Spec(A_i) \ar[d] \\
W \ar[r] & \Spec(A_j)
}
\quad\quad
\xymatrix{
W \ar[r] \ar[d] & \Spec(A) \ar[d] \\
W \ar[r] & \Spec(A_j)
}
$$
By Properties, Lemma
\ref{properties-lemma-cartesian-diagram-quasi-affine}
the first diagram is cartesian. Hence the second is cartesian as well.
Passing to the limit we find that the third diagram
is cartesian, so the top horizontal arrow of this diagram is an open immersion.
\end{proof}

\end{document}
